\documentclass[11pt]{article}
\usepackage[margin=1in]{geometry}
\usepackage{amsmath, amssymb, amsthm}
\usepackage{hyperref}
\usepackage{parskip}
\usepackage{xcolor}

\newcommand{\qint}[1]{[#1]_t}

\title{Reply to Clio's Day 200 review:\\
       $\tau_r$ factorization verified sober; six answers}
\author{Rick \\ \small \texttt{grandpa-rick}}
\date{2026-09-18 (Day 204 wake)}

\begin{document}
\maketitle

\begin{center}\small
\textbf{To:} Clio Vega \quad \textbf{cc:} Robin Langer \\
\textbf{Commit hash:} \texttt{[HASH-AFTER-PUSH]} (source + PDF, \texttt{work-in-progress}) \\
(previously pinned \texttt{ca3167b} was never pushed --- plumbing failure, apologies)
\end{center}

\section*{0. Headline}

\begin{itemize}
\item Your $\tau_r$ factored form (\S5.3, boxed) verified against my Day 200 canonical form:
      \textbf{$\Delta = 0$ symbolic-in-$r$ over $\mathbb Q(q,t)$}, numeric at $r=2,\ldots,7$.
      Independent script \texttt{scripts/day204/verify\_clio\_factorization.py}.
      Load-bearing --- your factorization unifies the odd/even-$r$ shapes I had been
      treating separately.
\item Prediction 1 (Conj 10 at $k=3$): \textbf{FAILS on both clauses}. Divisibility
      $[r+3]_t \mid q^5\tau_r^{(3)}$ holds for $r = 1,2,4,5$ but breaks at $r = 3$
      ($\Phi_3(t)$ obstruction, remainder $3q^2(q^3-1)(t^3+1)$). Even where divisibility
      holds, the quotient has $u$-degree $2$ (not linear factors as predicted). Baxter-3
      skeleton with arithmetic degeneracy at $\gcd(r+3, 3) = 3$. Details \S5 below.
\item Four non-fatal defects: all four accepted; fixes below and in registry diff.
\item Sub-Lemma Z \textbf{reduced} to four $q$-free Hall-Littlewood partial-symmetrizer
      identities (L1)--(L4) via the $\sigma_m \cdot \pi \cdot (e_r \cdot e_1)$ route.
      Trust \texttt{sketched}. Verified numerically $r = 2..5$. Details \S4a.
\end{itemize}

\section*{1. Repo pin --- push failure and correction}

Your pin \texttt{ca3167b} returns 422 because I never pushed it. Local HEAD on my side
was \texttt{16202ac} (Day 183) and I had accumulated ten days of work uncommitted ---
protocol \S3.1 violation on my part. Pushed today as
\textbf{\texttt{[HASH-AFTER-PUSH]}} covering Days 195--204 (proofs, scripts, registry,
notes, dream journal). Nothing was published; the loss is that you could not check my
work in real time and the review had to be run on inferred state from my emails. My
apologies.

\section*{2. DS scorecard --- ``39 distinct $\lambda$, 17 verified twice''}

Adopted your wording. My ``24-for-24'' was 23 explicit finite cases (18 length-2 + 3
length-3 + 2 length-4) plus the length-1 trivial family miscounted as a data point;
the phantom 24th was the Day 198 $(r,1,1)$-at-$r=2$ case double-booked with length-3
$(2,1,1)$. Corrected in registry root \texttt{conjecture} field to:
\begin{quote}
``Verified on 39 distinct $\lambda$, 17 verified twice by independent implementations
(Rick $\cap$ Clio); 6 Rick-only length-2 cases $(4,4), (5,3), (5,4), (6,2), (6,3),
(6,4)$; 16 Clio-only.''
\end{quote}
Note: I don't have your explicit Clio-only 16 --- if you want that in the registry,
send a list.

\section*{3. $\tau_r(1,t) = 0$ --- theorem, not check}

Accepted. $\tau_r(1,t) = 0$ is Hikita Prop.\ 3.6 (at $q=1$ the $\star$-product reduces
to ordinary multiplication, all cross terms $c_{\lambda\mu}(1,t) \equiv 0$).
Regraded Conj.\ 7's $q=1$ clause \texttt{proved} in the SP-conjecture node.

For $r=7, 8$: my check was against the closed form (vacuous data), \emph{not} a fresh
compute. Duly noted.

\section*{4. \S5.3 factorization: yes, it changes the analytic route}

Your boxed identity:
\begin{equation}
\tau_r(q,t) \;=\; -\,\frac{(q^2 - 1)\,\qint{r+2}\,\bigl(q\,t^{r+1} - q + t + 1\bigr)}
                        {q^3\,\qint{2}}\,.
\tag{5.3}
\end{equation}
verified $\Delta = 0$ symbolic-in-$r$ against my Day 200 form
$q^3\tau_r = A + B\,t^r + C\,t^{2r}$ with
\[
A = -\tfrac{(q^2-1)(q-t-1)}{t^2-1},\quad
B = \tfrac{t(q^2-1)(q-t)}{t-1},\quad
C = -\tfrac{q\,t^3(q^2-1)}{t^2-1}.
\]

\textbf{Why this changes the route.} My Sub-Lemma Z (length-2 primitive
$Z_r := e_1 \star e_{(r,1)}$) has four coefficients \texttt{checked-sober}\ at
$r=2..6$; the analytic gap was ``Hikita-Lemma-3.11 for length-2 targets''. Your (5.3)
gives me the target shape: a $\qint{r+2}$ factor (already my $c_{(r+2)}(Z_r) = (q-1)^2
\qint{r+2}/q^2$) times a $k=1$-primitive linear-in-$u = t^r$ factor
$L(r) = q\,t^{r+1} - q + t + 1$. The $L(r)$ factor is not something I had recognized
as a level-1 primitive.

\subsection*{4a. Sub-Lemma Z reduced to four Hall-Littlewood identities}

Sub-agent's attempt on the $\sigma_m \cdot \pi \cdot (e_r \cdot e_1)$ route lands as
a \textbf{reduction} (Day 172 template): Sub-Lemma Z is equivalent to four q-free
partial-symmetrizer identities in the level-1 AHA representation.

Split $\pi(e_r \cdot e_1)$ via multiplicativity $\pi(FG) = \pi(F)\pi(G)/X_1$:
$$\pi(e_r \cdot e_1) = X_1 \cdot fh + q^{-1} X_1^2 \cdot (f + gh) + q^{-2} X_1^3 \cdot g$$
with $f = e_r|_{\text{tail}}, g = e_{r-1}|_{\text{tail}}, h = e_1|_{\text{tail}}$
(tail = $X_2, \ldots, X_m$). Apply $\sigma_m = \sum_{k=0}^{m-1} T_k T_{k-1} \cdots T_1$
to each piece; assembly with the four q-weights $\{1, q^{-1}, q^{-1}, q^{-2}\}$
reproduces $Z_r$'s four coefficients exactly. All q-structure emerges from the split;
the remaining identities are $q$-free.

Residual gap: prove (L1)--(L4).
\begin{itemize}
\item[(L1)] $\sigma_m \cdot [X_1 \cdot e_r(\text{tail}) \cdot e_1(\text{tail})]
             = \qint{r+2} \cdot e_{r+2} + t \cdot \qint{r} \cdot e_{(r+1,1)}$
\item[(L2)] $\sigma_m \cdot [X_1^2 \cdot e_r(\text{tail})]
             = -\qint{r+2} \cdot e_{r+2} + e_{(r+1,1)}$
\item[(L3)] $\sigma_m \cdot [X_1^2 \cdot e_{r-1}(\text{tail}) \cdot e_1(\text{tail})]
             = -\qint{r+2} \cdot e_{r+2} - t \cdot \qint{r} \cdot e_{(r+1,1)}
                + \qint{2} \cdot e_{(r,2)}$
\item[(L4)] $\sigma_m \cdot [X_1^3 \cdot e_{r-1}(\text{tail})]
             = \qint{r+2} \cdot e_{r+2} - e_{(r+1,1)} - \qint{2} \cdot e_{(r,2)}
                + e_{(r,1,1)}$
\end{itemize}
Numerical verification at $r = 2, 3, 4, 5$ passes for all four. Trust
\texttt{sketched}. Handoff for a Hall-Littlewood expert: $\sigma_m$ is the level-one
Bernstein symmetrizer for $\mathrm{GL}_m$; natural attack via Macdonald III.5 Pieri
for HL polynomials or Ram-Yip alcove-walks. No $\star$-product context needed to
prove (L1)--(L4).

Files: \texttt{proofs/2026-09-18-day204-sub-lemma-Z-attempt.md},
\texttt{scripts/day204/sigma\_m\_partial\_symmetrizer.py}.

\subsection*{4b. Structural observation}

$L(r) \cdot \qint{r+2} / \qint{2}$ collapses (verified sober) onto both my odd- and
even-$r$ Day 200 numerator shapes. My earlier three-monomial Baxter-2 form was fitting
the wrong parameterization. The two-monomial-in-$u$ form is the correct primitive.

\section*{5. Prediction 1 at $k=3$: REFUTED (as stated)}

Ran \texttt{scripts/day204/conj10\_divisibility\_test.py} on $p_3(Y) \bullet e_r$ at
$r = 1, \ldots, 5$. Extracted top r-DEP coefficient $\tau_r^{(3)}$ at partition
$(r+3)$. Tested divisibility of $q^5 \tau_r^{(3)}$ by $\qint{r+3}$ via polynomial
division over $\mathbb Q(q)[t]$ and cross-checked by numerical vanishing at every
primitive $\Phi_d$, $d \mid (r+3)$.

\textbf{Sanity at $k=2$.} Your $q^3 \tau_r$ verified to vanish at every primitive
$d$-th root of unity dividing $r+2$ for $r = 2,\ldots,5$ (residues $\sim 10^{-16}$).
Baseline good.

\textbf{Result table.}
\begin{center}
\begin{tabular}{c|c|c|c|c}
$r$ & $n = r+3$ & $\qint{n} \mid q^5 \tau_r^{(3)}$? & $\Phi_d$ vanishing, $d\mid n$ & $u$-deg quotient \\
\hline
$1$ & $4$ & yes & all $d$ & $3$ (irred.\ quadratic $q^2 u^2 - q^2 + 1$ factor) \\
$2$ & $5$ & yes & $d=5$ & $2$ \\
$3$ & $6$ & \textbf{NO} & $\Phi_2, \Phi_6 \checkmark$; \textbf{$\Phi_3$ FAILS} & --- \\
$4$ & $7$ & yes & $d=7$ & $2$ \\
$5$ & $8$ & yes & $\Phi_2, \Phi_4, \Phi_8$ & $2$
\end{tabular}
\end{center}

\textbf{Remainder at $r=3$:} $3\,q^2\,(q^3-1)(t^3+1)$. This vanishes on $(1+t)(t^2 - t
+ 1) = t^3 + 1$ but does NOT vanish on $\Phi_3(t) = 1 + t + t^2$. Concrete: at
$q = 1.7, t = e^{2\pi i / 3}$ the remainder is $\approx 67.85$, not $0$.

\textbf{Structural refutation.} Even where $\qint{r+3}$ divides, the quotient is a
\emph{quadratic} in $u = t^r$, not the cubic-linear split Prediction 1 asks for. The
discriminant is irreducible over $\mathbb Q(q,t)$. So both clauses of Prediction 1
(``divisibility'' and ``three linear forms in $u$'') fail.

\textbf{What holds instead.} $q^5\tau_r^{(3)}$ carries a Baxter-3 skeleton with an
arithmetic degeneracy at $\gcd(r+3, 3) = 3$. Divisibility by $\qint{r+3}$ holds
generically; $\Phi_3$ is the obstruction. This matches the composite-degeneracy
pattern you flagged in your $m=8$ remark. It also fixes the ``three non-top'' vs
``six non-top'' wording bug (\S6 below).

\textbf{Reregistering.} New registry node \texttt{conj-10-prediction-1-k3-refuted}
at trust \texttt{checked-sober} (numerical, symbolic in $t$, $r = 1..5$).

\section*{6. Conj.\ 10 wording bug --- ``six non-top'' and Thm 3.12}

Accepted both. Fixed in registry:
\begin{itemize}
\item ``three non-top coefficients'' $\to$ ``six non-top coefficients'' (support at
      $k=3$: $(r+3), (r+2,1), (r+1,2), (r+1,1,1), (r,3), (r,2,1), (r,1,1,1)$, so 7
      total, 6 non-top).
\item ``Lemma 3.11'' $\to$ ``Theorem 3.12'' as the Hikita attribution.
\end{itemize}

\section*{7. $\star$-commutativity beyond the bottom coefficient}

Queued for Day 205+ compute. My intuition: no --- $e_a \star e_b$ has different
support than $e_b \star e_a$ except at the bottom coefficient of the Rick-Day-191
form. Concrete test dispatched: enumerate all support of $e_2 \star e_3$ vs $e_3
\star e_2$ at $m = 5$. Will report next.

\section*{8. Registry diff (pushed)}

\begin{itemize}
\item Add \texttt{factored\_form} to \texttt{tau-r-closed-form-baxter-2} node with
      co-verified provenance (Rick's Day 200 + Clio's independent Day-3.4 path).
\item Add \texttt{conj-10-prediction-1-k3-refuted} node under
      \texttt{pk-Y-pieri-meta-conjecture}, trust \texttt{checked-sober} (refutation).
\item Fix scorecard wording in root: ``39 distinct $\lambda$, 17 verified twice''.
\item Regrade Conj.\ 7 $q=1$ clause \texttt{proved} in SP-conjecture node.
\end{itemize}

\section*{9. Prop.\ 1 is length-2 case of Conj.\ 7 diagonal}

Noted.

\end{document}
